\documentclass{article}

\usepackage{amsmath,amssymb}
\usepackage{xurl}
\usepackage[hidelinks]{hyperref}
\setlength{\emergencystretch}{2em}

% Shared commands for notes in docs/paper-gaps/.

\DeclareUnicodeCharacter{2080}{\ifmmode{}_{0}\else\textsubscript{0}\fi}
\DeclareUnicodeCharacter{2081}{\ifmmode{}_{1}\else\textsubscript{1}\fi}
\DeclareUnicodeCharacter{2082}{\ifmmode{}_{2}\else\textsubscript{2}\fi}
\DeclareUnicodeCharacter{2099}{\ifmmode{}_{n}\else\textsubscript{n}\fi}

\newcommand{\Lean}{\textsc{Lean}}
\ExplSyntaxOn
\NewDocumentCommand{\leanid}{m}{
  \ifmmode
    \textsf{\detokenize{#1}}
  \else
    \group_begin:
    \urlstyle{sf}
    \tl_set:Nn \l_tmpa_tl {#1}
    \tl_replace_all:Nnn \l_tmpa_tl {\_} {_}
    \exp_args:NV \path \l_tmpa_tl
    \group_end:
  \fi
}
\ExplSyntaxOff
\providecommand{\tr}{\operatorname{tr}}
\newcommand{\tnleanIssueBase}{https://github.com/LionSR/TNLean/issues/}
\newcommand{\tnleanPrBase}{https://github.com/LionSR/TNLean/pull/}
\newcommand{\tnleanDocBase}{https://sirui-lu.com/TNLean/docs/}
\newcommand{\ghissue}[1]{\href{\tnleanIssueBase#1}{GitHub issue \##1}}
\newcommand{\ghpr}[1]{\href{\tnleanPrBase#1}{PR \##1}}
\newcommand{\doclink}[1]{\href{\tnleanDocBase#1.html}{\path{#1}}}

% Dirac notation and the identity operator, as in blueprint/src/macros/common.tex.
% \providecommand keeps note-local definitions authoritative.
\usepackage{dsfont}
\providecommand{\ket}[1]{|#1\rangle}
\providecommand{\bra}[1]{\langle #1|}
\providecommand{\braket}[2]{\langle #1 | #2 \rangle}
\providecommand{\ketbra}[2]{|#1\rangle\!\langle #2|}
\providecommand{\Id}{\mathds{1}}
\providecommand{\one}{\mathds{1}}
\providecommand{\C}{\mathbb{C}}
\providecommand{\MN}[1]{M_{#1}(\C)}
\providecommand{\MD}{M_D(\C)}

% External referencing for paper sources.  The build helper writes sanitized
% auxiliary files under build/paper-aux/ and excludes duplicated source labels.
\usepackage{xr}

% Import a generated paper auxiliary file with a caller-supplied label prefix.
% The candidate paths support builds from docs/paper-gaps/, the repository root,
% and build/paper-aux/ itself.
\newcommand{\importpaperaux}[2]{%
  \IfFileExists{../../build/paper-aux/#2.aux}{%
    \externaldocument[#1]{../../build/paper-aux/#2}%
  }{%
    \IfFileExists{build/paper-aux/#2.aux}{%
      \externaldocument[#1]{build/paper-aux/#2}%
    }{%
      \IfFileExists{#2.aux}{%
        \externaldocument[#1]{#2}%
      }{}%
    }%
  }%
}

% General external reference macros
\newcommand{\paperref}[2]{\ref{#1-#2}}
\newcommand{\papereqref}[2]{(\ref{#1-#2})}

% CPSV16 source (1606.00608 / MPDO-22-12-17-2.tex) external document and shortcuts
\importpaperaux{cpsv16-}{1606.00608/MPDO-22-12-17-2-tnlean}
\newcommand{\cpsvref}[1]{\paperref{cpsv16}{#1}}
\newcommand{\cpsveqref}[1]{\papereqref{cpsv16}{#1}}

% Verdict marker: \gapnote{<kind>}{<status>}. Kind is the policy
% classification (clarification, local-correction, scope-restriction,
% unfaithful, false-source, open-gap); status is open, wip (elimination
% underway), resolved, or historical. Machine-read by the site index;
% prints a verdict line.
\newcommand{\gapnote}[2]{\par\noindent\textbf{Verdict:} #1 (#2).\par}


\title{Common Blocking and Span Equality in the Non-periodic MPS Fundamental Theorem}
\author{}
\date{2026-04-30}

\begin{document}
\maketitle
\gapnote{clarification}{resolved}

\section{The compressed source argument}

The non-periodic matrix product state fundamental theorem reduces arbitrary
tensors with the same matrix product vector families to comparison data for
bases of normal tensors. The relevant sources are the MPDO canonical-form paper
and its published version
\cite{Cirac2016MPDO_arXiv,Cirac2017MPDO_AnnPhys}, together with the later review
\cite{Cirac2021Matrix}.\footnote{For traceability, this note corresponds to
\href{https://github.com/LionSR/TNLean/issues/1047}{issue \#1047}.}
The cited arguments are not in question. This note separates the canonical-form
steps that the papers compress but the formal proof must state explicitly.

Let $A=(A^i)_{i=1}^d$ be a tensor with physical dimension $d$. For every
positive integer $N$, its matrix product vector is
\begin{align}
    \ket{V^{(N)}(A)}
        &=
        \sum_{i_1,\ldots,i_N=1}^{d}
        \tr\left(A^{i_1}\cdots A^{i_N}\right)
        \ket{i_1}\otimes\cdots\otimes\ket{i_N}.
    \label{eq:rmp_common_blocking_span_equality_mpv_definition}
\end{align}
This is the positive-length matrix product vector family defined in
Ref.~\cite[lines~130--138]{Cirac2016MPDO_arXiv}.

After removing invariant subspaces, the source places the tensor in the block
form
\begin{align}
    A^i
        &=
        \bigoplus_{k=1}^{r}\mu_k A_k^i.
    \label{eq:rmp_common_blocking_span_equality_canonical_blocks}
\end{align}
Zero blocks are allowed, while every nonzero block is scaled so that its
transfer map has spectral radius one. A nonzero irreducible block may retain a
nontrivial peripheral period. Blocking by a common multiple of all such periods
removes this periodicity
\cite[lines~214--231]{Cirac2016MPDO_arXiv}. The review uses the same
canonical-form reduction \cite[lines~1798--1820]{Cirac2021Matrix}.

The subsequent comparison uses bases of normal tensors. A basis of normal
tensors for $A$ is a minimal family $(A_j)_{j=1}^{g}$ of normal tensors such
that every vector in the matrix product vector family of $A$ is a linear
combination of the corresponding vectors of the $A_j$, which are eventually
linearly independent. Equivalently, every normal block in the canonical form is
a phase times a gauge transform of exactly one basis tensor, and the basis
family is minimal
\cite[Proposition~\texttt{prop:char-BNT}, lines~271--280]{Cirac2016MPDO_arXiv}.
The review gives the same characterization
\cite[Proposition~\texttt{prop:char-BNT}, lines~1846--1859]{Cirac2021Matrix}.

If the blocks associated with $A_j$ have weights
$\mu_{j,1},\ldots,\mu_{j,r_j}$, then the canonical-form tensor satisfies
\begin{align}
    \ket{V^{(N)}(A)}
        &=
        \sum_{j=1}^{g}
        \left(\sum_{q=1}^{r_j}\mu_{j,q}^{N}\right)
        \ket{V^{(N)}(A_j)}.
    \label{eq:rmp_common_blocking_span_equality_bnt_expansion}
\end{align}
The equal-case proof uses this expansion in
Ref.~\cite[lines~283--302]{Cirac2016MPDO_arXiv}; the review uses the same
formula \cite[lines~1864--1885]{Cirac2021Matrix}.

Injectivity enters only after this canonical-form construction. The source
defines injective and block-injective canonical forms and invokes the quantum
Wielandt theorem to show that every normal tensor becomes injective after a
bounded additional blocking
\cite[lines~317--332]{Cirac2016MPDO_arXiv}. The review gives the corresponding
$3D^5$ bound for obtaining block-injective canonical form
\cite[lines~2120--2124]{Cirac2021Matrix}. Period removal and injectivity are
therefore distinct operations: period removal produces primitive normal blocks,
whereas later overlap-rigidity arguments require one-site injectivity or an
explicit substitute obtained after further blocking.

The source treats the equal-case comparison compactly. Equality or
proportionality of the two matrix product vector families first matches their
bases of normal tensors up to phases and gauges. After relabeling the matched
basis tensors, the coefficient comparison becomes
\begin{align}
    \sum_{q=1}^{r_{a,j}}\mu_{j,q}^{N}
        &=
        \sum_{q=1}^{r_{b,j}}
        \left(\nu_{j,q}e^{i\phi_j}\right)^{N}
        \qquad
        \text{for every $j$ and $N$}.
    \label{eq:rmp_common_blocking_span_equality_power_sum_comparison}
\end{align}
An elementary power-sum argument then recovers the multiplicities and weights,
after which the blockwise gauges assemble into a global conjugating matrix
\cite[lines~1181--1192]{Cirac2016MPDO_arXiv}. The review states the resulting
proportional and equal matrix product vector fundamental theorems in the same
form \cite[lines~1891--1900]{Cirac2021Matrix}.

\section{The formal argument}

The canonical-form reduction uses the positive-length relation
\leanid{SameMPV₂Pos}, which matches the source definition in
(\ref{eq:rmp_common_blocking_span_equality_mpv_definition}). All-zero blocks
contribute nothing at positive length and are discarded. The structural
direct-sum identity directly bounds the total bond dimension of the retained
blocks, without using an empty-word comparison.\footnote{See
\leanid{exists_irreducible_blockDecomp_nonzeroBlocks} and
\leanid{exists_tp_gauge_from_arbitrary}.}

The common blocking step must also be explicit. For one block family, the source
instruction to ``block by the least common multiple of the periods'' is
unambiguous. For two tensors, however, the formal proof must choose one positive
physical blocking length for both families. If the period-removal length of an
original block is $m_k$, the common length is written as $p=m_ke_k$ for a
positive integer $e_k$.

Blocking first by $m_k$ and then by $e_k$ produces physical labels consisting
of nested words. Blocking once by $p$ instead produces labels in the alphabet
of words of length $p$. The formal statement must therefore retain the
relabeling between these two alphabets.\footnote{The common-sector data and the
relabeling results are recorded in
\path{blueprint/src/chapter/ch09_canonical.tex:1645--1824} and
\path{blueprint/src/chapter/ch09_canonical.tex:3230--3284}. The main formal
statement is
\texttt{MPSTensor.fundamentalTheorem\_after\_blocking\_commonLength\_commonSector};
see
\path{TNLean/MPS/CanonicalForm/SectorComparison/CommonSectorData.lean} and
\path{TNLean/MPS/CanonicalForm/SectorComparison/CommonSectorTransport.lean}.}
Consequently, the canonical-form theorem supplies relabeled common-sector
families and a conditional label-compatibility statement. It does not
immediately identify every blocked canonical nonzero part with its flattened
sector tensor.

The downstream formal comparison now follows the coefficient-comparison route
for basis-of-normal-tensor sectors. Its data consist of the matched-sector phase
relation, coefficient identities obtained from eventual linear independence,
copy-weight matching obtained from power sums, and a block-diagonal global
gauge. Further Wielandt blocking is needed only by arguments that require
one-site injectivity; it is not part of the current coefficient-comparison
theorem.\footnote{The final comparison route is recorded in
\path{blueprint/src/chapter/ch11_fundamental_theorem_proof.tex:31--460}.}

Let $P_j$ and $Q_k$ denote matched basis-of-normal-tensor sectors, let
$\beta(k)$ be the matching permutation, let $\zeta_k$ be the associated phase,
and let $\sigma$ denote a physical word. The matched-sector relation is
\begin{align}
    V^{(N)}(Q_k)_{\sigma}
        &=
        \zeta_k^{N}
        V^{(N)}\left(P_{\beta(k)}\right)_{\sigma}.
    \label{eq:rmp_common_blocking_span_equality_sector_phase_relation}
\end{align}
If $c_N^{(j)}(P)$ and $c_N^{(k)}(Q)$ denote the corresponding sector
coefficients, eventual linear independence gives
\begin{align}
    c_N^{(\beta(k))}(P)
        &=
        \zeta_k^{N}c_N^{(k)}(Q).
    \label{eq:rmp_common_blocking_span_equality_sector_coefficient_identity}
\end{align}
Newton--Girard power-sum comparison then identifies the copy weights, and the
matched block gauges assemble into the direct-sum global gauge.\footnote{See
\path{blueprint/src/chapter/ch11_fundamental_theorem_proof.tex:31--460},
especially the sector coefficient identity, copy-weight matching, and assembled
global-gauge entries. The corresponding formal objects live in
\path{TNLean/MPS/FundamentalTheorem/SectorBNT/CoeffIdentity.lean},
\path{TNLean/MPS/FundamentalTheorem/SectorBNT/CopyWeightMatching.lean},
\path{TNLean/MPS/FundamentalTheorem/SectorBNT/Fundamental.lean}, and
\path{TNLean/MPS/FundamentalTheorem/SectorBNT/FundamentalCoord.lean}. The
removed comparison lemmas from the former span-based route are no longer used
in the final comparison chapter; the remaining proof follows this
coefficient-comparison route.}
The coefficient-identity and copy-weight lemmas state the exact comparison data
required downstream. In the proportional case, the proof still requires the
coefficient comparison for the phases supplied by the sector-matching theorem.

\section{Why the distinctions belong together}

The conversion from equality of arbitrary matrix product vector families to the
hypotheses of the basis-of-normal-tensor sector comparison couples these
canonical-form operations. Common blocking determines both the physical
alphabet and the relabeled sector families to which the coefficient comparison
applies. Any later injectivity refinement belongs to a separate argument, such
as overlap rigidity, and is not a hypothesis of the SectorBNT coefficient
route. The matched-sector coefficient identities and copy-weight data refer to
the relabeled nonzero-sector families, not to a family before the
common-alphabet relabeling.

The formal proof therefore needs the complete chain: discard all-zero blocks at
positive lengths, choose one common positive blocking length, transport through
the word relabeling, keep later overlap-rigidity hypotheses separate, and
convert basis-of-normal-tensor matching into the coefficient and copy-weight
identities used by the global gauge comparison. The blueprint presents this
chain as the connection between the canonical-form chapter and the final
comparison chapter.\footnote{See
\path{blueprint/src/chapter/ch11_fundamental_theorem_proof.tex:539--629} for the
current equal and proportional BNT canonical-form statements. The proportional
global gauge remains conditional on the coefficient identity recorded at
\path{blueprint/src/chapter/ch11_fundamental_theorem_proof.tex:361--388}; the
equal case uses
\path{blueprint/src/chapter/ch11_fundamental_theorem_proof.tex:503--537}.}

The same-side equal-norm scalar obstruction from the broader
fundamental-theorem development is not a flaw in the CPSV or CPGSV arguments.
That obstruction corrected an over-strong formal statement that identified an
equal-norm class with one basis tensor. The source expansion
(\ref{eq:rmp_common_blocking_span_equality_bnt_expansion}) already permits
multiple basis tensors and multiplicities. The formal development therefore
replaced naive norm-class grouping by phase-class sector decompositions and
common-cover data.

\section{Conclusion}

The cited papers compress several standard canonical-form operations, and this
note gives no counterexample. Zero or nilpotent contributions are invisible to
positive-length matrix product vectors; blocking removes peripheral periods; a
later Wielandt-type blocking yields injectivity; phase and gauge data match the
bases of normal tensors; and power sums recover coefficient multiplicities.
The formal development separates these operations because it transports the
matrix product vector relation through explicit positive blocking, places both
tensors over one explicitly relabeled alphabet, and constructs the matched
sector coefficient identities and copy-weight equations only after that
relabeling.

The discrepancy is therefore neither a counterexample nor a theorem-level
source error. It explains why the formal proof cannot replace the full
canonical-form route by one invocation of basis-of-normal-tensor uniqueness.
The required comparison points are removal of all-zero blocks, construction of
common primitive blocks, common blocking, iterated-blocking relabeling,
matched-sector coefficient identities, and copy-weight matching.\footnote{For
traceability, these points are recorded in
\href{https://github.com/LionSR/TNLean/issues/652}{issue \#652},
\href{https://github.com/LionSR/TNLean/issues/942}{issue \#942},
\href{https://github.com/LionSR/TNLean/issues/969}{issue \#969},
\href{https://github.com/LionSR/TNLean/issues/970}{issue \#970},
\href{https://github.com/LionSR/TNLean/issues/971}{issue \#971},
\href{https://github.com/LionSR/TNLean/issues/944}{issue \#944}, and
\href{https://github.com/LionSR/TNLean/issues/990}{issue \#990}.}

\bibliographystyle{alpha}
\bibliography{references}

\end{document}
